\section{Hypothesis}\label{P2}
The statistical method known as hypothesis testing is typically utilized 
in the field of scientific research in order to arrive at a conclusion regarding 
the parameters of a data set (such as the mean, variance, and so on).
\par
On the basis of two competing hypotheses, the hypothesis for a particular situation 
might be described as follows: The null hypothesis (H0) asserts that the two procedures 
that are being compared are of identical quality 
(in accordance with the evaluation methods that will be utilized subsequently). 
On the other hand, the alternative hypothesis, often known as Ha, is something that is 
defined as the assertion that one approach is superior to the other method. 
Both of these hypotheses are false if one of them is true, and vice versa if the other 
one becomes true. The rejection of the null hypothesis is one way that can be used to 
demonstrate that the alternative hypothesis is supported by evidence and that one method 
is superior to the other method. On the basis of this, the null and alternative hypotheses 
should be defined in the following manner as the initial stage in the process of assessing 
the two codes contained in this report:

\subsection{Null Hypothesis}
Both the results produced from my code and the results got from the code of  my colleague 
have the same performance capabilities and are comparable to one another. 
The average amount of time required to complete this work is the same for both algorithms 
(i.e., $\mu_A$ = $\mu_B$), which means that we are unable to tell which algorithm is more efficient 
than the other during the process of gathering the boxes.

\subsection{Alternative hypothesis}
When it comes to grouping the boxes next to each other under a variety of conditions, 
the algorithm written by my colleague is far faster than my own approach. Consequently, 
the average amount of time that my robot takes to collect all of the boxes is longer 
than the amount of time that my colleague's code takes (i.e., $\mu_A$ is greater than $\mu_B$).

\subsection{Tests}
The first thing that has to be done in order to test the null hypothesis is to collect 
some sample test data. The amount of data collected throughout the exam ought to 
be sufficient in order to produce results that are reliable, accurate, and appropriate. 
The Z-test, the T-test, the two sampled T-test, the paired T-test, and other similar tests 
are some of the various approaches that can be utilized to examine and compare the performance 
of two algorithms. In order for us to be able to proceed with the analysis using a certain 
approach, there are a few fundamental preconditions that must be met. Typically, 
the Z-test is utilized in situations when the population, as well as its mean and 
standard deviation, are accessible. When using the T-test, it is necessary to have knowledge 
about at least one of the characteristics of the population, such as the mean or the 
standard deviation. For the purpose of this research, the population is not available; 
rather, we are merely contrasting two distinct algorithms that are capable of completing a task. 
The paired T-test and the two sampled T-test are both potential usefulnesses for this analysis.

\newpage